\documentclass[aspectratio=169,10pt]{beamer}
\usetheme{default}
\setbeamertemplate{navigation symbols}{}
\usepackage{graphicx}\usepackage[T1]{fontenc}\usepackage{amsmath,amssymb}
\usepackage{listings}
\definecolor{ncsured}{HTML}{CC0000}\definecolor{ncsudk}{HTML}{990000}
\definecolor{ncsugray}{HTML}{6D6D6D}\definecolor{ncsuteal}{HTML}{008473}
\definecolor{codebg}{HTML}{F5F5F5}
\setbeamercolor{structure}{fg=ncsured}
\setbeamercolor{frametitle}{fg=white,bg=ncsured}
\setbeamercolor{title}{fg=ncsudk}
\setbeamercolor{itemize item}{fg=ncsured}
\setbeamercolor{itemize subitem}{fg=ncsured}
\setbeamerfont{frametitle}{series=\bfseries}
\setbeamertemplate{frametitle}{\nointerlineskip\vskip2pt%
  \begin{beamercolorbox}[wd=\paperwidth,ht=2.4ex,dp=1ex,leftskip=.3cm]{frametitle}%
  \small\bfseries\insertframetitle\end{beamercolorbox}}
\setbeamertemplate{footline}{\hbox{\begin{beamercolorbox}[wd=\paperwidth,ht=2.2ex,dp=1ex,leftskip=.3cm,rightskip=.3cm]{}%
  \tiny Optimal Market Making for Cryptocurrency Options \hfill Week 4 \hfill \insertframenumber\end{beamercolorbox}}}
\lstset{basicstyle=\ttfamily\scriptsize,backgroundcolor=\color{codebg},
  keywordstyle=\color{ncsured},commentstyle=\color{ncsugray},
  stringstyle=\color{ncsuteal},breaklines=true,columns=fullflexible,
  showstringspaces=false,frame=single,rulecolor=\color{ncsugray}}

\title{Week 4: Stochastic Optimal Control}
\subtitle{The general theory behind the Avellaneda-Stoikov worked example}
\date{North Carolina State University \,\textbar\, Summer 2026}
\titlegraphic{
\includegraphics[width=4.2cm]{/Users/dendisuhubdy/Github/ncsu/assets/ncsu_logo.png}}
\begin{document}
\begin{frame}[plain]\titlepage\begin{center}\tiny Industry Mentor: Dendi Suhubdy (Bitwyre)  |  Guest: Fenni Kang (Coincall)\end{center}\end{frame}
\begin{frame}{Learning Objectives}
\begin{itemize}
  \item Apply the dynamic programming principle in continuous time
  \item Understand the HJB equation and the verification theorem
  \item See why CARA / CRRA utilities give tractable problems
\end{itemize}
\end{frame}

\begin{frame}{Lecture 1: Dynamic Programming \& HJB}
\begin{itemize}
  \item Controlled diffusions and admissible controls
  \item The dynamic programming principle (Bellman)
  \item The infinitesimal generator and the HJB equation
\end{itemize}
\end{frame}

\begin{frame}{Lecture 2: Verification \& Utilities}
\begin{itemize}
  \item Verification theorem: a smooth HJB solution IS the value function
  \item CARA vs CRRA: where wealth enters the optimal policy
  \item Why market making uses CARA
\end{itemize}
\end{frame}

\begin{frame}{Illustration}
\begin{center}\includegraphics[width=0.86\textwidth]{/Users/dendisuhubdy/Github/ncsu/weekly_presentations/figures/week04.png}\end{center}
\end{frame}

\section*{Derivation}
\begin{frame}{Derivation \& Rationale}
\[V(t,x)=\sup_{u}\ \mathbb{E}\!\left[\int_t^T f(s,X_s,u_s)\,ds + g(X_T)\right]\]

\small\textbf{\color{ncsured}Why:} The value function is the best achievable expected reward from $(t,x)$ onward. We define the whole problem through $V$ because Bellman's principle then gives a local (differential) equation instead of an intractable global optimization.
\end{frame}

\begin{frame}{Derivation \& Rationale}
\[V(t,x)=\sup_u \mathbb{E}\big[V(t{+}dt, X_{t+dt}) + f\,dt\big]\]

\small\textbf{\color{ncsured}Why:} The dynamic programming principle: an optimal policy is optimal from every later state too. Expanding $V(t+dt,\cdot)$ with Ito's lemma turns this into the HJB PDE.
\end{frame}

\begin{frame}{Derivation \& Rationale}
\[\partial_t V + \sup_{u}\big\{\mathcal{L}^u V + f(t,x,u)\big\}=0,\quad V(T,x)=g(x)\]

\small\textbf{\color{ncsured}Why:} The HJB equation. $\mathcal{L}^u$ is the generator of the controlled diffusion; the $\sup$ is the instantaneous optimal control. This is the engine behind every model in the course -- Avellaneda-Stoikov is one special case.
\end{frame}

\begin{frame}{Derivation \& Rationale}
\[\text{CARA: } U(w)=-e^{-\gamma w}\ \Rightarrow\ \text{policy independent of } w\]

\small\textbf{\color{ncsured}Why:} Constant absolute risk aversion makes marginal utility scale-free in wealth, so the optimal quotes do not move with the cash balance. CRRA ($U=w^{1-\gamma}/(1-\gamma)$) instead makes the policy scale WITH wealth -- natural for investment, awkward for a desk.
\end{frame}

\begin{frame}{Getting the Data from Coincall}
\begin{itemize}
  \item Theory week -- no Coincall data required.
  \item Optional: estimate drift/vol from perp klines to parameterize Merton.
\end{itemize}
\vspace{2pt}\tiny\textit{Cross-check exact paths at https://docs.coincall.com/}
\end{frame}

\begin{frame}[fragile]{Code --- week4\_merton.py}
\begin{lstlisting}
import torch
# Merton under CARA, solved by backward value iteration on a wealth grid.
w = torch.linspace(-5, 5, 401); dt = 1e-3; mu, r, sig, gam = 0.08, 0.02, 0.2, 1.0
V = -torch.exp(-gam * w)                      # terminal utility
for _ in range(1000):                         # march backward in time
    # optimal risky holding pi* = (mu - r)/(gam*sig^2) is wealth-independent (CARA)
    pi = (mu - r) / (gam * sig**2)
    drift = (r*w + pi*(mu - r))
    diff  = 0.5 * (pi*sig)**2
    dV = torch.gradient(V, spacing=(w[1]-w[0]).item())[0]
    d2V = torch.gradient(dV, spacing=(w[1]-w[0]).item())[0]
    V = V + dt * (drift*dV + diff*d2V)
print("pi* (constant in wealth):", (mu - r)/(gam*sig**2))
\end{lstlisting}
\end{frame}
\begin{frame}{Glossary --- Key Terms}
\footnotesize\begin{description}
\item[\color{ncsured}Stochastic optimal control] choosing controls to optimize an expected objective under randomness
\item[\color{ncsured}Dynamic programming principle] an optimal policy stays optimal from every later state
\item[\color{ncsured}Value function] the best achievable expected reward from a given state onward
\item[\color{ncsured}Infinitesimal generator] operator giving the expected instantaneous change of a function
\item[\color{ncsured}Verification theorem] conditions under which an HJB solution is the value function
\item[\color{ncsured}CRRA utility] constant relative risk aversion; the policy scales with wealth
\item[\color{ncsured}Admissible control] a control satisfying the problem's constraints
\end{description}
\end{frame}

\begin{frame}{Reading}
\begin{itemize}
  \item Pham (2009), Ch. 3-4; Cartea-Jaikumar-Penalva, Ch. 6
\end{itemize}
\end{frame}

\begin{frame}{References}
\scriptsize\begin{itemize}\setlength{\itemsep}{3pt}
  \item Pham, H. (2009). Continuous-time Stochastic Control and Optimization with Financial Applications. Springer.
  \item Fleming, W. H., \& Soner, H. M. (2006). Controlled Markov Processes and Viscosity Solutions (2nd ed.). Springer.
  \item Oksendal, B. (2003). Stochastic Differential Equations (6th ed.). Springer.
  \item Almgren, R., \& Chriss, N. (2001). Optimal execution of portfolio transactions. Journal of Risk, 3(2), 5-40.
\end{itemize}
\end{frame}

\end{document}